\documentclass[12pt, a4paper, reqno, oneside]{book}
\usepackage[left=2.5cm, right=2cm, top=2cm, bottom=2cm]{geometry}
\usepackage{amsmath,amsxtra,amssymb,latexsym, amscd,amsthm}
\usepackage{indentfirst}
\usepackage{color}
\usepackage{empheq}
\usepackage[USenglish]{babel}
% \usepackage[nocaptions]{vntex}
\usepackage{amsmath}
\usepackage[mathscr]{eucal}
\usepackage{amsfonts}
\usepackage{graphicx}
\usepackage{subcaption}
\usepackage{float}
\usepackage{xcolor}
\usepackage{verbatim}
%\usepackage{tikz}
\usepackage{tkz-graph}
\usetikzlibrary {positioning}
\usepackage{float}
\usepackage{array}
\usepackage{multirow}
%\usepackage{algorithm}
\usepackage[ruled,vlined,linesnumbered, algochapter]{algorithm2e}
% \usepackage{algpseudocode} 
\usepackage{pseudocode} % Chau sua de compile tren Linux
\usepackage{titlesec}
\usepackage{fancyhdr}
\fancyhf{}
\fancyhead[LE,RO]{}%\small{\rightmark}
\fancyhead[LO,RE]{\small{\leftmark}}
\fancyfoot[C]{ \thepage}
\usepackage{ulem}
\usepackage{pdfpages}
\usepackage{niceframe}
\usepackage[square,numbers]{natbib}
\usepackage{subcaption,booktabs}

%\input setbmp
\usepackage[hyphens,spaces,obeyspaces]{url}
\usepackage[unicode,breaklinks=true]{hyperref}
% \usepackage{url}

% \usepackage{enumitem}
%\usepackage{titledot}
%\titlename{Chương}
%+ formating
%\setlength{\oddsidemargin}{0in}
%\setlength{\textwidth}{6.in}
%\setlength{\topmargin}{-0.5in}
%\setlength{\textheight}{9.25in}
\renewcommand{\baselinestretch}{1.25}

%Chương không số
\newcommand \chap[1]{%
  \chapter*{#1}%
  \addcontentsline{toc}{chapter}{#1}
  \markboth{#1}{}
}
\newcommand \sect[1]{%
  \section*{#1}%
  \addcontentsline{toc}{section}{#1}}
  
%+ Dinh nghia mot so ky hieu toan hoc.
\def\B{\mathscr B}
\def\one{\mathbf 1}
\def\R{{\mathbb R}}
\def\LL{{\mathbb L}}
\newcommand\abs[1]{|#1|}
\newcommand\set[1]{\{#1\}}
\newcommand\norm[1]{||#1||}
\DeclareMathOperator{\vrai}{vrai}
\DeclareMathOperator{\const}{const}
\newcommand{\eoproof}{\hfill $\square$}
\newcommand{\spro}[1]{\left<#1\right>}
\newcommand{\argmin}{\arg\min}

\renewcommand{\theequation}{\arabic{chapter}.\arabic{equation}} %\arabic{section}.
\newtheorem{lemma}{\bf Bổ đề}[chapter]
\newtheorem{theorem}{\bf Định lý}[chapter]
\newtheorem{proposition}{\bf Mệnh đề}[chapter]
\newtheorem{corollary}{\bf Hệ quả}[chapter]
\newtheorem{definition}{\bf Định nghĩa}[chapter]
\newtheorem{remark}{\bf Chú ý}[section]
% \newtheorem{problem}{Bài toán}
\renewcommand*{\listalgorithmcfname}{Danh sách thuật toán}
\renewcommand*{\algorithmcfname}{Thuật toán}
\renewcommand*{\algorithmautorefname}{thuật toán}

%mức tiêu đề
\titleclass{\subsubsubsection}{straight}[\subsection]

\newcounter{subsubsubsection}[subsubsection]
\renewcommand\thesubsubsubsection{\thesubsubsection.\arabic{subsubsubsection}}
\renewcommand\theparagraph{\thesubsubsubsection.\arabic{paragraph}} % optional; useful if paragraphs are to be numbered

\titleformat{\subsubsubsection}
  {\normalfont\normalsize\bfseries}{\thesubsubsubsection}{1em}{}
\titlespacing*{\subsubsubsection}
{0pt}{3.25ex plus 1ex minus .2ex}{1.5ex plus .2ex}

\makeatletter
\renewcommand\paragraph{\@startsection{paragraph}{5}{\z@}%
  {3.25ex \@plus1ex \@minus.2ex}%
  {-1em}%
  {\normalfont\normalsize\bfseries}}
\renewcommand\subparagraph{\@startsection{subparagraph}{6}{\parindent}%
  {3.25ex \@plus1ex \@minus .2ex}%
  {-1em}%
  {\normalfont\normalsize\bfseries}}
\def\toclevel@subsubsubsection{4}
\def\toclevel@paragraph{5}
\def\toclevel@paragraph{6}
\def\l@subsubsubsection{\@dottedtocline{4}{7em}{4em}}
\def\l@paragraph{\@dottedtocline{5}{10em}{5em}}
\def\l@subparagraph{\@dottedtocline{6}{14em}{6em}}
\makeatother

\setcounter{secnumdepth}{4}
\setcounter{tocdepth}{4}


%thư mục hình
\graphicspath{{./figs/}}

\definecolor{mygreen}{rgb}{0,0.6,0}
\definecolor{mygray}{rgb}{0.5,0.5,0.5}
\definecolor{mymauve}{rgb}{0.58,0,0.82}

\lstdefinestyle{customc}{
  belowcaptionskip=1\baselineskip,
  breaklines=true,
  frame=single,
  xleftmargin=\parindent,
  language=C,
  showstringspaces=false,
  basicstyle=\footnotesize\ttfamily,
  keywordstyle=\bfseries\color{green!40!black},
  commentstyle=\itshape\color{purple!40!black},
  identifierstyle=\color{blue},
  stringstyle=\color{orange},
  numbers=left,                   % where to put the line-numbers
  numberstyle=\tiny\color{gray},  % the style that is used for the line-numbers
  stepnumber=1, 
}

\lstdefinestyle{customasm}{
  belowcaptionskip=1\baselineskip,
  frame=single,
  xleftmargin=\parindent,
  language=[x86masm]Assembler,
  basicstyle=\footnotesize\ttfamily,
  commentstyle=\itshape\color{purple!40!black},
  numbers=left,                   % where to put the line-numbers
  numberstyle=\tiny\color{gray},  % the style that is used for the line-numbers
  stepnumber=1, 
}

\lstset{
  language=Python,
  basicstyle=\ttfamily\small,
  keywordstyle=\color{blue},
  commentstyle=\color{green},
  stringstyle=\color{red},
  numbers=left,
  numberstyle=\tiny\color{gray},
  stepnumber=1,
  numbersep=10pt,
  frame=single,
  breaklines=true,
  showspaces=false,
  showstringspaces=false,
  showtabs=false,
  tabsize=2
}

\definecolor {processblue}{cmyk}{0.96,0,0,0}

\newcolumntype{C}[1]{>{\centering\arraybackslash}p{#1}}

% \algnewcommand\algorithmicrequir{\textbf{Require:}}
% \algnewcommand\REQUIRE{\item[\algorithmicrequir]}

% \algnewcommand\algorithmicinput{\textbf{Input:}}
% \algnewcommand\INPUT{\item[\algorithmicinput]}
% \algnewcommand\algorithmicoutput{\textbf{Output:}}
% \algnewcommand\OUTPUT{\item[\algorithmicoutput]}

% \makeatletter
% \renewcommand{\ALG@beginalgorithmic}{\scriptsize}
% \makeatother

\SetKwProg{KwFn}{Function}{}{} 